% =========================================================
% The Four-Layer Hierarchy
% =========================================================
\paragraph{The Four-Layer Hierarchy.}

The twelve moves are not all at the same level. They form a four-layer
hierarchy that reflects the logical structure of metric space proofs:
each layer depends on the one below it, and together the four layers
describe a progression from the most basic (unpacking definitions) to
the most structural (global topological reasoning).

Understanding the hierarchy is not an organizational convenience.
It tells you something deep: \emph{why} metric spaces are defined the
way they are, and why nearly every proof in the subject follows the
same basic upward shape. Once you see this, metric space theory stops
looking like a collection of unrelated tricks and starts looking like
a small, elegant machine built from a few core mechanisms.

\begin{tcolorbox}[colback=gray!6, colframe=gray!40, arc=2pt,
  left=6pt, right=6pt, top=4pt, bottom=4pt,
  title={\small\textbf{The Four Levels}},
  fonttitle=\small\bfseries]
\small
\renewcommand{\arraystretch}{1.4}
\begin{tabular}{p{0.06\textwidth}p{0.22\textwidth}p{0.30\textwidth}p{0.30\textwidth}}
\toprule
\textbf{L} & \textbf{Name} & \textbf{Moves} & \textbf{What it does} \\
\midrule
L1 & Definition Expansion & Move 1 &
  Turns abstract concepts into inequalities \\
L2 & Inequality Mechanics & Moves 2, 3, 5, 9 &
  Combines and bounds distances \\
L3 & Sequence Control & Moves 4, 6, 8, 12 &
  Controls asymptotic and tail behavior \\
L4 & Structural / Topological & Moves 7, 10, 11 &
  Reasons about the space as a whole \\
\bottomrule
\end{tabular}
\end{tcolorbox}

\begin{remark}[Level 1: Definition Expansion --- the foundation of everything]
Every other move operates on inequalities. Those inequalities do not
appear until you unpack a definition. Before Level~1, you have a
statement about abstract objects. After Level~1, you have quantifiers
and specific inequality conditions that can be manipulated.

This is why Level~1 is always the first step in any analysis proof,
regardless of how complex the argument becomes. The inequality
$d(x_n, p) < \varepsilon$ does not exist in the proof until you write
it down. You write it down by unpacking the definition of convergence.

What Level~1 produces:
\begin{itemize}[itemsep=2pt]
  \item ``$x_n \to p$'' $\mapsto$ $\forall\,\varepsilon>0\;\exists\,N$
    s.t.\ $n \geq N \Rightarrow d(x_n, p) < \varepsilon$
  \item ``$\{x_n\}$ is Cauchy'' $\mapsto$ $\forall\,\varepsilon>0\;\exists\,M$
    s.t.\ $m,n \geq M \Rightarrow d(x_m, x_n) < \varepsilon$
  \item ``$K$ is compact'' $\mapsto$ every open cover of $K$ has a finite subcover
  \item ``$E$ is closed'' $\mapsto$ every convergent sequence in $E$ converges
    to a point of $E$
\end{itemize}
The abstract word is a compressed form of the inequality. Unpacking
it is not simplification; it is \emph{activation}.
\end{remark}

\begin{remark}[Level 2: Inequality Mechanics --- the algebra of metric spaces]
Once Level~1 has produced inequalities, Level~2 manipulates them.
This is the algebra of metric spaces: you combine, bound, split, and
chain inequalities to move from what you know toward what you need.

The triangle inequality is a \emph{metric axiom} for a reason: it was
included in the definition of a metric precisely because it is the tool
needed to chain distance bounds. Without it, you cannot get from control
over $d(x,y)$ and $d(y,z)$ to control over $d(x,z)$. The axiom is there
because it enables Level~2.

Notice that Level~2 does not involve sequences, indices, or tails.
It is purely about bounding distances between specific points. The
sequence machinery comes at the next level.
\end{remark}

\begin{remark}[Level 3: Sequence Control --- the engine of limit arguments]
Analysis is fundamentally about sequences approaching limits, distances
becoming arbitrarily small, behavior stabilizing as an index grows large.
Level~3 controls this asymptotic behavior.

The key idea: sequence control is about managing the word ``eventually.''
The definition of convergence says that for large enough $n$, the
distance $d(x_n, p)$ is small. The tail argument combines multiple
``large enough $n$'' conditions by taking the maximum index. Choosing $N$
explicitly produces the specific threshold. Subsequence extraction selects
which terms of the sequence to keep.

These moves require genuine sequence machinery --- you are reasoning
about indexed infinite processes, not just about individual distances
between points. That is what places them at Level~3 rather than Level~2.
\end{remark}

\begin{remark}[Level 4: Structural / Topological --- global reasoning]
Level~4 moves operate on the global structure of the metric space, not
on individual sequences or inequalities. Contradiction reasons about
the logical structure of the whole proof. Embedding reasons about how
the space sits inside a larger ambient space. Compactness reasons about
the global covering properties of the space.

Compactness is the flagship Level~4 property. It converts infinite
problems into finite ones: an infinite family of open sets reduces to
a finite subcover, and a finite collection has a minimum, a maximum, a
single controlling constant. This is the most powerful tool in metric
space theory.

Level~4 and Level~3 interact bidirectionally. Sequential compactness
(every sequence has a convergent subsequence) is a Level~3
characterization of a Level~4 property. Many deep results in analysis
are precisely about these interactions.
\end{remark}

\begin{remark}[Reading a proof using the hierarchy]
When you read a proof in the textbook, label each step by level.
This is not a mechanical exercise; it reveals the logical architecture
and tells you why each step appears where it does.

Here is the proof that convergent implies Cauchy, labeled:

\begin{quote}
Let $\varepsilon > 0$.\hfill \textit{[L1: introduce $\varepsilon$ from the Cauchy definition]}

Since $x_n \to p$, there exists $N$ such that $n \geq N \Rightarrow d(x_n,p) < \varepsilon/2$.
\hfill \textit{[L3: choose $N$ from the convergence hypothesis]}

For $m, n \geq N$:
\hfill \textit{[L3: tail condition]}

$d(x_m, x_n) \leq d(x_m, p) + d(p, x_n)$
\hfill \textit{[L2: triangle inequality]}

$\phantom{d(x_m, x_n)} < \varepsilon/2 + \varepsilon/2 = \varepsilon$.\hfill
\textit{[L2: $\varepsilon$-splitting]}

Hence $\{x_n\}$ is Cauchy.\hfill \textit{[L1: conclusion matches Cauchy definition]}
\end{quote}

This proof uses Levels 1, 2, and 3, but not Level 4. It is a local
argument about sequences, not a global structural argument. To reach
Level~4, you would need compactness or an embedding argument.

The levels a proof reaches tell you its depth and what tools it uses.
A verification exercise stays at Levels~1--2. A limit argument reaches
Level~3. A compactness argument reaches Level~4.
\end{remark}

\begin{remark}[Why the metric axioms are the axioms they are]
The metric axioms --- non-negativity, identity of indiscernibles,
symmetry, triangle inequality --- are not arbitrary. They are exactly
the conditions needed to make Level~2 possible.

Without the triangle inequality, you cannot chain distance bounds.
Without identity of indiscernibles, the notion of a limit point
degenerates. Without symmetry, the definitions of balls and convergence
become asymmetric and unworkable.

The axioms are the \emph{minimal conditions on a distance function
that support the full four-level proof hierarchy.} When you study a
new property of metric spaces, ask: which level does it live at?
The answer tells you what tools to expect and what kind of proofs
will be required.
\end{remark}
